% Zumdahl Chapter 8 -- Bonding: General Concepts. ELABORATED notes, 5 blocks.
\documentclass{apchem-notes}

\apcunitword{Chapter}
\apcunit{8}
\apcunittitle{Bonding: General Concepts}
\apcdoctitle{Guided Notes}
\apcsubtitle{Zumdahl \S8.1--8.13 \textbullet\ PDF pp.~390--453 \textbullet\ 5 blocks}

\begin{document}
\apctitleblock

\begin{apcnote}[How this chapter fits the AP course]
A high-value chapter: nearly all of \textbf{CED Unit 2} lives here, and it
also supplies two topics that belong to other units ---
\textbf{\S8.4} (ions and ionic radii) is CED 1.8, and
\textbf{\S8.8} (bond energies) is CED 6.7, the bond-enthalpy topic Unit 6
borrows from this chapter.

Blocks 4 and 5 cover Lewis structures, resonance, and VSEPR, which you
worked in Unit 2. Treat those as consolidation; the genuinely new material
for most students is \textbf{lattice energy} (Block 2) and
\textbf{computing $\Delta H$ from bond energies} (Block 3).

\S8.6 (partial ionic character) is background reading only.
\end{apcnote}

% =====================================================================
\blocklesson{Bond Types, Electronegativity, and Polarity}{Zumdahl \S8.1--8.3}

\begin{objectives}
  \item Classify bonds using electronegativity difference.
  \item Distinguish bond polarity from molecular polarity.
\end{objectives}

\reading{Zumdahl \S8.1--8.3}{391--400}

\begin{warmup}[3]
  \item Which is more electronegative, N or O?
        \answerblank[2cm]{O}
  \item Geometry of \ce{CO2}: \answerblank[2.4cm]{linear}
  \item Configuration of \ce{Mg^2+}: \answerblank[2.6cm]{$[\ce{Ne}]$}
\end{warmup}

\segment{INSTRUCTION A}{25}{Electronegativity}

\notesectionz{The property that decides bond type}{8.2}
\practice{6}

\term{Electronegativity} is an atom's ability to attract shared electrons in
a bond. On the Pauling scale it runs from
\answerblank[1.8cm]{4.0} for fluorine down to \answerblank[1.8cm]{0.7} for
caesium.

The trend follows the same Coulombic engine as everything in Unit 1: it
\answerblank[2cm]{increases} across a period (rising $Z_{\text{eff}}$) and
\answerblank[2cm]{decreases} down a group (greater distance and shielding).

\renewcommand{\arraystretch}{1.4}
\noindent\begin{tabular}{@{}p{3.4cm}p{3.2cm}p{6.2cm}@{}}
\toprule
\textbf{Bond type} & \textbf{$\Delta$EN} & \textbf{Electron behaviour}\\
\midrule
Nonpolar covalent & \answerblank[1.4cm]{$\approx 0$} &
  shared equally; no partial charges\\
Polar covalent & intermediate &
  shared \answerblank[2.2cm]{unequally}; $\delta^+$ and $\delta^-$ appear\\
Ionic & large &
  \answerblank[2.4cm]{transferred} to form ions\\
\bottomrule
\end{tabular}

\begin{apcnote}
These are regions on a \emph{continuum}, not three separate boxes. Zumdahl
is explicit about this: as $\Delta$EN grows, bonds acquire progressively
more ionic character. Do not argue about where the boundary sits --- argue
with the $\Delta$EN value.
\end{apcnote}

\begin{workedexample}[Worked example: ranking bond polarity]
Order \ce{H-H}, \ce{S-H}, \ce{Cl-H}, \ce{O-H}, \ce{F-H} by polarity.

Electronegativities: H 2.1, S 2.5, Cl 3.0, O 3.5, F 4.0. So
$\Delta$EN $=$ 0, 0.4, 0.9, 1.4, 1.9 respectively:
\[ \ce{H-H} < \ce{S-H} < \ce{Cl-H} < \ce{O-H} < \ce{F-H} \]
\end{workedexample}

\segment{GUIDED PRACTICE}{15}{Classify and rank}

\begin{enumerate}[leftmargin=*,itemsep=0.3em]
  \item \ce{K-Br}: \answerblank[5cm]{ionic --- large $\Delta$EN}
  \item \ce{N-N}: \answerblank[5cm]{nonpolar covalent}
  \item \ce{C-O}: \answerblank[5cm]{polar covalent}
  \item Most polar of \ce{C-N}, \ce{C-O}, \ce{C-F}:
        \answerblank[3cm]{\ce{C-F}}
\end{enumerate}

\segment{INSTRUCTION B}{20}{From bond dipoles to molecular dipoles}

\notesectionz{Dipole moments}{8.3}
\practice{1}

A polar bond has a \term{bond dipole} --- a vector pointing toward the more
electronegative atom. Whether the \emph{molecule} is polar depends on
whether those vectors \answerblank[1.8cm]{cancel}.

Two conditions, both required for a polar molecule:
\begin{enumerate}[leftmargin=*,itemsep=0.15em]
  \item at least one \answerblank[2.4cm]{polar bond}, and
  \item a geometry that does \emph{not} produce
        \answerblank[2.6cm]{cancellation}.
\end{enumerate}

\begin{apctrap}
``It has polar bonds, so it is polar'' is the single most common error on
this topic. \ce{CO2} has two strongly polar C=O bonds and is
\textbf{nonpolar}, because its linear geometry points them exactly opposite.
Same story for \ce{BF3} (trigonal planar) and \ce{CCl4} (tetrahedral).
\end{apctrap}

\segment{APPLICATION}{20}{Polarity decisions}

\begin{enumerate}[leftmargin=*,itemsep=0.4em]
  \item \ce{CCl4} and \ce{CHCl3} both contain polar C--Cl bonds, yet only
        one is a polar molecule. Identify it and explain.
        \answerlines[3]{\ce{CHCl3}. In \ce{CCl4} four identical bond dipoles
        arranged tetrahedrally sum to zero. Replacing one Cl with H breaks
        that symmetry, so the dipoles no longer cancel and a net dipole
        remains.}
  \item \ce{NH3} is polar but \ce{BF3} is not, even though both are
        ``\ce{XY3}.'' Explain in terms of geometry.
        \answerlines[3]{\ce{NH3} has a lone pair on nitrogen, giving four
        electron domains and a trigonal pyramidal shape whose dipoles do not
        cancel. \ce{BF3} has only three domains and is trigonal planar, so
        its three equal dipoles at 120$^\circ$ sum to zero.}
\end{enumerate}

\begin{exitticket}
Can a molecule containing only nonpolar bonds be polar? Explain.
\answerlines[2]{No --- with no bond dipoles there is nothing to fail to
cancel, so the molecular dipole must be zero.}
\end{exitticket}

% =====================================================================
\blocklesson{Ions and Lattice Energy}{Zumdahl \S8.4--8.5}

\begin{objectives}
  \item Predict ion configurations and rank ionic radii, including
        isoelectronic series.
  \item Use Coulomb's law to rank lattice energies.
\end{objectives}

\reading{Zumdahl \S8.4--8.5}{400--408}

\begin{warmup}[3]
  \item $\Delta$EN of \ce{H-F} given H 2.1, F 4.0:
        \answerblank[2cm]{1.9}
  \item Is \ce{BF3} polar? \answerblank[2cm]{no}
  \item Configuration of \ce{S^2-}: \answerblank[2.6cm]{$[\ce{Ar}]$}
\end{warmup}

\segment{INSTRUCTION A}{25}{Ion sizes}

\notesectionz{How ions differ from their parent atoms}{8.4}
\practice{6}

\begin{itemize}[leftmargin=*,itemsep=0.25em]
  \item A \term{cation} is always \answerblank[2cm]{smaller} than its parent
        atom --- electrons are removed (often emptying a whole shell) while
        the nuclear charge is unchanged, so the remainder is pulled in
        tightly.
  \item An \term{anion} is always \answerblank[2cm]{larger} --- added
        electrons increase mutual repulsion with no extra protons to
        compensate.
\end{itemize}

\notesub{Isoelectronic series}

Ions with the same electron count differ only in nuclear charge. Zumdahl's
rule: for a series of isoelectronic ions, size
\answerblank[2.2cm]{decreases} as $Z$ increases.

\begin{workedexample}[Worked example: Zumdahl's krypton series]
Rank \ce{Se^2-}, \ce{Br-}, \ce{Rb+}, \ce{Sr^2+} by decreasing size.

All four have the krypton configuration --- 36 electrons each --- so
repulsions are comparable. Their nuclear charges are $Z = 34, 35, 37, 38$.
More protons pulling the same 36 electrons means a smaller ion:
\[ \ce{Se^2-} > \ce{Br-} > \ce{Rb+} > \ce{Sr^2+} \]
\end{workedexample}

\segment{GUIDED PRACTICE}{15}{Rank the ions}

\begin{enumerate}[leftmargin=*,itemsep=0.3em]
  \item Larger: \ce{Cl} or \ce{Cl-}? \answerblank[2.4cm]{\ce{Cl-}}
  \item Larger: \ce{Fe^2+} or \ce{Fe^3+}? \answerblank[2.4cm]{\ce{Fe^2+}}
  \item Rank by size: \ce{O^2-}, \ce{F-}, \ce{Na+}, \ce{Mg^2+}
        \answerblank[6cm]{\ce{O^2-} $>$ \ce{F-} $>$ \ce{Na+} $>$ \ce{Mg^2+}}
\end{enumerate}

\segment{INSTRUCTION B}{20}{Lattice energy}

\notesectionz{Coulomb's law, applied to a crystal}{8.5}
\practice{5}

\term{Lattice energy} is the energy change when gaseous ions assemble into
an ionic solid. Zumdahl writes it as a modified Coulomb's law:
\[
  \text{Lattice energy} = k\!\left(\answerblank[1.8cm]{$\dfrac{Q_1Q_2}{r}$}\right)
\]
where $Q_1$ and $Q_2$ are the ionic \answerblank[1.6cm]{charges} and $r$ is
the distance between ion \answerblank[1.6cm]{centres}.

\notesub{The ranking procedure}

\begin{enumerate}[leftmargin=*,itemsep=0.2em]
  \item Compare the \answerblank[3.4cm]{charge products} $|Q_1Q_2|$ first
        --- they dominate.
  \item Only if those tie, compare \answerblank[2cm]{ionic radii}: smaller
        ions sit closer, giving a stronger lattice.
\end{enumerate}

\begin{workedexample}[Worked example: three-way ranking]
Rank \ce{NaCl}, \ce{MgO}, and \ce{NaF} by lattice energy magnitude.

Charge products: \ce{MgO} is $(2)(2) = 4$; \ce{NaCl} and \ce{NaF} are both
$(1)(1) = 1$. So \ce{MgO} is far strongest.

Tiebreak \ce{NaF} vs \ce{NaCl}: \ce{F-} is smaller than \ce{Cl-}, so $r$ is
shorter and the attraction stronger.
\[ \ce{MgO} \gg \ce{NaF} > \ce{NaCl} \]
\end{workedexample}

\notesub{Why this matters structurally}

In the \ce{LiF} lattice each \ce{Li+} is surrounded by
\answerblank[1.2cm]{six} \ce{F-} ions and vice versa --- the sodium chloride
structure. That packing maximizes attractions between opposite charges while
minimizing repulsions between like charges.

\segment{APPLICATION}{20}{Lattice energy reasoning}

\begin{enumerate}[leftmargin=*,itemsep=0.4em]
  \item \ce{MgO} melts at \SI{2852}{\celsius} but \ce{NaF} at
        \SI{993}{\celsius}, despite similar ion sizes. Explain
        quantitatively. \answerlines[3]{Charge product is 4 for \ce{MgO}
        versus 1 for \ce{NaF}, so the Coulombic attraction is roughly four
        times as strong at comparable separation --- far more energy is
        needed to disrupt the lattice.}
  \item Rank \ce{CaO}, \ce{KCl}, \ce{KBr}, \ce{SrO} by lattice energy and
        justify. \answerlines[3]{\ce{CaO} $>$ \ce{SrO} $\gg$ \ce{KCl} $>$
        \ce{KBr}. Charges first: the oxides are $2{+}/2{-}$. Within each
        pair, the smaller ion wins --- \ce{Ca^2+} $<$ \ce{Sr^2+} and
        \ce{Cl-} $<$ \ce{Br-}.}
\end{enumerate}

\begin{exitticket}
Why is \ce{Al2O3} used as a furnace lining while \ce{NaCl} is not?
\answerlines[2]{\ce{Al2O3} has a $3{+}/2{-}$ charge product of 6 versus 1
for \ce{NaCl}, so its lattice energy and melting point are enormously
higher.}
\end{exitticket}

% =====================================================================
\blocklesson{Bond Energies}{Zumdahl \S8.7--8.8}

\begin{objectives}
  \item Relate bond order to bond length and bond energy.
  \item Estimate $\Delta H$ for a reaction from bond energies.
\end{objectives}

\reading{Zumdahl \S8.7--8.8}{409--415}

\begin{warmup}[4]
  \item Lattice energy depends on charges and what else?
        \answerblank[2.6cm]{ionic radii}
  \item Larger: \ce{K+} or \ce{Ca^2+}? \answerblank[2.4cm]{\ce{K+}}
  \item Stronger lattice: \ce{LiF} or \ce{KI}?
        \answerblank[2.4cm]{\ce{LiF}}
  \item Sign of $\Delta H$ for an exothermic reaction:
        \answerblank[2.2cm]{negative}
\end{warmup}

\segment{INSTRUCTION A}{25}{What a bond energy is}

\notesectionz{Breaking bonds always costs energy}{8.8}
\practice{5}

The \term{bond energy} is the energy required to break one mole of a
particular bond in the gas phase. It is always
\answerblank[1.8cm]{positive} --- breaking a bond is always
\answerblank[2.7cm]{endothermic}, and forming one always releases energy.

\renewcommand{\arraystretch}{1.3}
\noindent\begin{tabular}{@{}lccc@{}}
\toprule
 & \textbf{Bond order} & \textbf{Length} & \textbf{Energy (kJ/mol)}\\
\midrule
\ce{C-C} & 1 & longest & 347\\
\ce{C=C} & 2 & middle & 614\\
\ce{C#C} & 3 & \answerblank[1.6cm]{shortest} & \answerblank[1.4cm]{839}\\
\bottomrule
\end{tabular}

Higher bond order $\to$ \answerblank[1.8cm]{shorter} and
\answerblank[1.8cm]{stronger}. Note that the triple bond is not three times
the single --- the increments shrink.

\segment{GUIDED PRACTICE}{15}{Bond order reasoning}

\begin{enumerate}[leftmargin=*,itemsep=0.3em]
  \item Shortest N--O bond: \ce{NO3-}, \ce{NO2-}, or \ce{NO+}?
        \answerblank[3.2cm]{\ce{NO+} (order 3)}
  \item Which requires more energy to break, \ce{O2} or \ce{N2}? Why?
        \answerblank[5.5cm]{\ce{N2} --- triple bond}
\end{enumerate}

\segment{INSTRUCTION B}{20}{Computing $\Delta H$ from bonds}

\notesectionz{The bookkeeping}{8.8}
\practice{5}

Imagine the reaction proceeding in two stages: tear every reactant apart
into atoms, then assemble the products.
\[
  \Delta H = \answerblank[8.2cm]{$\sum D(\text{bonds broken}) - \sum D(\text{bonds formed})$}
\]

\begin{apctrap}
The order is \textbf{broken minus formed} --- the reverse of the formation
enthalpy equation from Chapter 6, which is products minus reactants.
Reversing them flips your sign. A quick sanity check: if the products have
stronger bonds, the reaction should come out exothermic.
\end{apctrap}

\begin{workedexample}[Worked example 1: Zumdahl's HF synthesis]
\[ \ce{H2(g) + F2(g) -> 2 HF(g)} \]
Bonds broken: one \ce{H-H} (432) and one \ce{F-F} (154).
Bonds formed: two \ce{H-F} (565 each).
\[
  \Delta H = (432 + 154) - 2(565) = 586 - 1130 = \SI{-544}{\kilo\joule}
\]
Compare with formation data: $2 \times (-271) = \SI{-542}{\kilo\joule}$.
The two methods agree to within 2 kJ.
\end{workedexample}

\begin{workedexample}[Worked example 2: HCl synthesis]
\[ \ce{H2(g) + Cl2(g) -> 2 HCl(g)} \]
Broken: \ce{H-H} 432 $+$ \ce{Cl-Cl} 239 $= 671$.
Formed: $2 \times \ce{H-Cl}$ 427 $= 854$.
\[ \Delta H = 671 - 854 = \SI{-183}{\kilo\joule} \]
Formation data give $2(-92.3) = \SI{-184.6}{\kilo\joule}$ --- again close.
\end{workedexample}

\begin{apcnote}
Bond energies give \emph{estimates}, not exact values, because tabulated
figures are averages over many compounds. A C--H bond in methane is not
identical to one in ethanol. Expect agreement within a few percent, and say
so if a question asks why the two methods differ.
\end{apcnote}

\segment{APPLICATION}{20}{Full calculation}

Bond energies (kJ/mol): \ce{N#N} 941, \ce{H-H} 432, \ce{N-H} 391.

\begin{enumerate}[leftmargin=*,itemsep=0.4em]
  \item Compute $\Delta H$ for \ce{N2(g) + 3 H2(g) -> 2 NH3(g)}.
        \workspace[2.6cm]
        \keyonly{\begin{apcanswer}Broken: $941 + 3(432) = 2237$. Formed:
        $6(391) = 2346$. $\Delta H = 2237 - 2346 =
        \SI{-109}{\kilo\joule}$\end{apcanswer}}
  \item The accepted value is \SI{-92.2}{\kilo\joule}. Explain the
        discrepancy. \answerlines[2]{Tabulated bond energies are averages
        across many molecules, so a calculation using them yields an
        estimate rather than an exact value.}
  \item Which contributed more to the exothermicity: breaking the triple
        bond, or forming the six N--H bonds? Explain.
        \answerlines[2]{Forming the N--H bonds --- that step releases
        2346 kJ, more than the 2237 kJ required to break all the reactant
        bonds, which is exactly why the reaction is exothermic.}
\end{enumerate}

\begin{exitticket}
Why is breaking a bond always endothermic? \answerlines[2]{A bond exists
because the bonded state is lower in energy than separated atoms; climbing
back out of that energy well requires an input of energy.}
\end{exitticket}

% =====================================================================
\blocklesson{Lewis Structures and Octet Exceptions}{Zumdahl \S8.9--8.11}

\begin{objectives}
  \item Draw Lewis structures reliably.
  \item Recognize the three classes of octet exception.
\end{objectives}

\reading{Zumdahl \S8.9--8.11}{415--423}

\begin{warmup}[3]
  \item $\Delta H$ formula from bond energies:
        \answerblank[4.6cm]{broken $-$ formed}
  \item Valence electrons in \ce{SO4^2-}:
        \answerblank[2.9cm]{$6+24+2 = 32$}
  \item Bond order of \ce{N2}: \answerblank[1.6cm]{3}
\end{warmup}

\segment{INSTRUCTION A}{25}{The procedure}

\notesectionz{Five steps}{8.10}
\practice{1}

\begin{enumerate}[leftmargin=*,itemsep=0.2em]
  \item Count total valence electrons ---
        \answerblank[7.1cm]{add for anions, subtract for cations}.
  \item Choose the central atom: the
        \answerblank[4.2cm]{least electronegative} one, never H.
  \item Connect with single bonds; complete the
        \answerblank[2.6cm]{outer atoms'} octets (H takes 2).
  \item Place any remaining electrons on the
        \answerblank[2.8cm]{central atom}.
  \item If the central atom is short, convert lone pairs into
        \answerblank[3.2cm]{multiple bonds}.
\end{enumerate}

Always finish with a two-part check: the total electron count \emph{and}
every octet.

\begin{workedexample}[Worked example: \ce{HCN}]
Count: $1 + 4 + 5 = 10$ electrons. Carbon is least electronegative (H can
never be central), so the skeleton is H--C--N.

Single bonds use 4; putting the remaining 6 on N gives N an octet but leaves
C with only 4. Convert two of nitrogen's lone pairs into bonds:
\[ \ce{H-C#N} \]
with one lone pair left on N. Check: 10 electrons, C and N both have
octets, H has 2. \checkmark
\end{workedexample}

\segment{GUIDED PRACTICE}{15}{Draw four}

Give the electron count, then the structure: \ce{PH3}, \ce{CO2},
\ce{SO2}, \ce{ClO4-}.
\workspace[4.4cm]
\keyonly{\begin{apcanswer}\ce{PH3}: 8 e$^-$, P central, 1 lone pair.
\ce{CO2}: 16 e$^-$, O=C=O, 2 lone pairs per O. \ce{SO2}: 18 e$^-$, bent,
lone pair on S. \ce{ClO4-}: 32 e$^-$, Cl central, 4 single bonds, 3 lone
pairs per O, brackets with $-$.\end{apcanswer}}

\segment{INSTRUCTION B}{20}{When the octet rule fails}

\notesectionz{Three classes of exception}{8.11}
\practice{1}

\renewcommand{\arraystretch}{1.35}
\noindent\begin{tabular}{@{}p{3.4cm}p{3.6cm}p{5.8cm}@{}}
\toprule
\textbf{Class} & \textbf{Examples} & \textbf{Why}\\
\midrule
Odd number of electrons & \ce{NO}, \ce{NO2} &
  one electron must remain \answerblank[1.8cm]{unpaired}\\
Fewer than an octet & \ce{BF3}, \ce{BeH2} &
  B and Be are content with \answerblank[1.6cm]{6} and 4\\
Expanded octet & \ce{SF6}, \ce{PCl5}, \ce{XeF2} &
  central atom is in period \answerblank[1.4cm]{3} or beyond\\
\bottomrule
\end{tabular}

\begin{apctrap}
Only period 3 and beyond can expand. Nitrogen \emph{cannot} hold more than
eight electrons, so structures like ``\ce{NF5}'' do not exist. If a
structure needs an expanded octet on a period-2 atom, the structure is
wrong.
\end{apctrap}

\segment{APPLICATION}{20}{Diagnose the structures}

\begin{enumerate}[leftmargin=*,itemsep=0.4em]
  \item A student draws \ce{BF3} with one B=F double bond to give boron an
        octet. Evaluate. \answerlines[3]{Rejected. The double bond forces a
        $+1$ formal charge onto fluorine, the most electronegative element
        --- energetically unreasonable. Boron's six-electron structure is
        accepted.}
  \item Explain why \ce{SF6} can exist but \ce{OF6} cannot.
        \answerlines[2]{Sulfur is in period 3 and can accommodate an
        expanded octet; oxygen is in period 2 and is limited to eight
        electrons.}
\end{enumerate}

\begin{exitticket}
Give the electron count for \ce{PO4^3-} and state whether an expanded octet
is required. \answerlines[2]{$5 + 24 + 3 = 32$ electrons. Not required ---
four single bonds and octets on the oxygens use exactly 32, though phosphorus
could expand if a different structure were drawn.}
\end{exitticket}

% =====================================================================
\blocklesson{Resonance, Formal Charge, and VSEPR}{Zumdahl \S8.12--8.13}

\begin{objectives}
  \item Use formal charge to select the best structure.
  \item Draw resonance sets and state their physical meaning.
  \item Predict geometry and bond angles with VSEPR.
\end{objectives}

\reading{Zumdahl \S8.12--8.13}{423--440}

\begin{warmup}[4]
  \item Electrons in \ce{NO3-}: \answerblank[2.4cm]{24}
  \item Can nitrogen expand its octet? \answerblank[1.6cm]{no}
  \item Which class of exception is \ce{BF3}?
        \answerblank[4.1cm]{fewer than an octet}
  \item Central atom of \ce{OCS}: \answerblank[1.8cm]{C}
\end{warmup}

\segment{INSTRUCTION A}{25}{Formal charge and resonance}

\notesectionz{Choosing among candidate structures}{8.12}
\practice{5}

\[
  FC = \text{valence} - \answerblank[5cm]{$\text{nonbonding} - \tfrac12(\text{bonding})$}
\]

The best structure has formal charges closest to
\answerblank[1.2cm]{zero}, with any negative charge on the
\answerblank[4.3cm]{most electronegative} atom. The formal charges must sum
to the species' overall charge --- that is a check, not a preference.

\notesub{Resonance}

When the electrons could be arranged in more than one equivalent way, draw
all forms joined by \answerblank[4.5cm]{double-headed arrows}. The real
molecule is the \term{hybrid} --- a single averaged structure, not a
molecule flipping between forms.

\begin{workedexample}[Worked example: nitrate]
\ce{NO3-} has 24 electrons and three equivalent forms, with the double bond
on each oxygen in turn.

The evidence that this is an average rather than an alternation: all three
N--O bonds are measured to be \emph{identical} in length, intermediate
between a single and a double bond. Bond order is
$4/3$, and each oxygen carries an average charge of $-2/3$.
\end{workedexample}

\begin{apctrap}
``The double bond moves back and forth'' scores zero. Cite the experimental
evidence instead: equal bond lengths, fractional bond order, delocalized
charge.
\end{apctrap}

\segment{GUIDED PRACTICE}{15}{Formal charge practice}

\begin{enumerate}[leftmargin=*,itemsep=0.3em]
  \item FC on S in \ce{SO4^2-} drawn with four single bonds:
        \answerblank[2cm]{$+2$}
  \item FC on each O in that structure: \answerblank[2cm]{$-1$}
  \item Bond order of C--O in \ce{CO3^2-}: \answerblank[2cm]{$4/3$}
  \item Bond order of O--O in \ce{O3}: \answerblank[2cm]{$3/2$}
\end{enumerate}

\segment{INSTRUCTION B}{20}{VSEPR}

\notesectionz{Geometry from electron-domain repulsion}{8.13}
\practice{1}

Electron domains --- bonding groups plus lone pairs --- arrange themselves
to maximize separation. A multiple bond counts as
\answerblank[2.4cm]{one domain}.

\renewcommand{\arraystretch}{1.3}
\noindent\begin{tabular}{@{}cllc@{}}
\toprule
\textbf{Domains} & \textbf{Domain geometry} & \textbf{Angle} & \textbf{Example}\\
\midrule
2 & \answerblank[2.6cm]{linear} & 180$^\circ$ & \ce{CO2}\\
3 & \answerblank[3cm]{trigonal planar} & 120$^\circ$ & \ce{BF3}\\
4 & \answerblank[2.6cm]{tetrahedral} & 109.5$^\circ$ & \ce{CH4}\\
5 & trigonal bipyramidal & 90/120$^\circ$ & \ce{PCl5}\\
6 & octahedral & 90$^\circ$ & \ce{SF6}\\
\bottomrule
\end{tabular}

Molecular geometry names only the \emph{atom} positions, so lone pairs
change the name: 4 domains with 1 lone pair is
\answerblank[3.9cm]{trigonal pyramidal}; with 2 lone pairs it is
\answerblank[1.6cm]{bent}.

Lone pairs also compress angles, because they repel
\answerblank[3.0cm]{more strongly} than bonding pairs --- hence
\ce{CH4} 109.5$^\circ$, \ce{NH3} $\approx107^\circ$, \ce{H2O}
$\approx104.5^\circ$.

\segment{APPLICATION}{20}{The full chain}

\begin{enumerate}[leftmargin=*,itemsep=0.4em]
  \item For \ce{SO2}: electron count, domains, geometry, angle, polarity.
        \answerlines[3]{18 e$^-$; 3 domains including one lone pair on S;
        bent; slightly less than 120$^\circ$; polar, because the bent shape
        prevents the S--O dipoles from cancelling.}
  \item \ce{CO2} and \ce{SO2} are both triatomic oxides, yet one is linear
        and one is bent. Explain the difference completely.
        \answerlines[3]{Carbon has two domains (two double bonds, no lone
        pairs) so \ce{CO2} is linear at 180$^\circ$. Sulfur has three
        domains including a lone pair, which occupies space and pushes the
        oxygens together into a bent shape.}
\end{enumerate}

\begin{exitticket}
Predict the geometry and approximate bond angle of \ce{NF3}, and state
whether it is polar. \answerlines[2]{Four domains with one lone pair $\to$
trigonal pyramidal, about 107$^\circ$. Polar --- the three N--F dipoles do
not cancel in that geometry.}
\end{exitticket}

\end{document}
